\minitopic{“人工在环”必须拥有实际能力}若审核者只看到模型结论，没有原始材料；每项任务只有数秒；其专业训练与任务不匹配；即使发现错误也无权阻止执行，那么一次点击确认不能构成有效监督。真实人机监督至少需要信息、时间、能力和否决权。四项中任何一项缺失，都应把责任重新分配给系统设计者与部署机构，而不能以“最终由人批准”吸收制度缺陷。

\minitopic{自动化偏差与算法厌恶是对称风险}用户可能因系统通常准确而忽略反证，也可能因一次显眼错误拒绝总体表现较好的工具。合理依赖不是固定地信或不信，而要按任务、亚群、环境和错误类型校准。界面应显示系统在哪些条件下验证、哪些输入超出分布、不同错误的历史频率以及人类覆写表现。只显示总体准确率会把分布性失败藏在平均数后，也让用户无法学习何时应依赖。

\minitopic{平均性能不能回答公正问题}面部分类系统在不同交叉群体上的误差差异说明，总体指标可能掩盖集中伤害\parencite{buolamwini2018}。即使两个系统总准确率相同，误报与漏报落在谁身上、后果是否对称、群体是否有申诉渠道，都可能不同。公平指标之间还会冲突，不能用一张分组表自动得出唯一结论。部署者必须说明采用何种公平目标、为什么适合任务、冲突如何处理以及哪些群体数据不足以支持判断。

\minitopic{责任图比“谁按了按钮”更有解释力}训练数据提供者影响表示范围，模型团队选择目标和损失函数，产品团队设计默认值与警告，机构决定部署场景与资源，使用者作出个案判断，管理者又决定申诉和停用条件。一次伤害可能来自多层共同作用。责任图应记录每层可预见什么、能控制什么、留下什么证据；控制能力与受益程度通常比因果链末端更能说明责任。这样既避免把工具拟人化为唯一责任者，也避免让复杂系统成为无人负责的借口。

\begin{argumentbox}
高风险人机系统上线前应有四类门槛：证据门槛——目标人群与任务上的验证；程序门槛——复核、覆写和升级路径；分配门槛——分组误差与后果评估；可修正门槛——版本记录、事故报告、停用条件和申诉。四类门槛缺一，平均准确率再高也不足以证明可部署。
\end{argumentbox}

\minitopic{事故处理是认识制度的一部分}系统错误后只修正单个输出，会失去学习机会。完整事件记录应保存输入、模型版本、检索材料、用户操作、当时警告、覆写理由和实际后果；随后判断是数据漂移、界面诱导、模型推断、权限不足还是政策设置造成。修复也要有层级：更正个案、通知受影响者、更新模型或流程、重新评估相似历史案例。若机构不知道错误发生过，就无法校准对系统的信任。

\minitopic{机器是否知道取决于理论与归属对象}可靠主义可能强调系统稳定产生真输出，能力论要求成功源于系统相关能力，责任主义则关注主体能否理解理由和响应批评。不同理论会对“机器本身知道”给出不同答案。更重要的实践问题常是人机整体是否产生可使用知识：哪个组件接触证据，哪个组件整合，谁能识别击败者，谁承担断言和行动。把归属对象明确后，争论才能避免在拟人语言与纯工具语言之间摆动。

\minitopic{标准框架提供检查表而非自动担保}风险管理框架可以要求映射用途、测量风险、治理责任和持续监测，例如NIST生成式AI风险管理资料所组织的任务\parencite{nist2024}。遵循框架不等于系统已经安全，更不等于每项输出为真。它的认识价值在于让假设、责任与证据要求公开，使不同部署能够被比较。框架若沦为勾选表，也会产生“合规即可靠”的新自动化偏差。

\minitopic{部署前评价必须逼近真实任务}通用基准上的高分不等于在医院、课堂或公共服务中可靠。任务语言、输入质量、用户专业水平与错误后果都会改变表现；测试集若与训练材料重叠，还会把记忆伪装成能力。部署者应建立本地案例集，包含常见任务、边缘条件、对抗输入与受影响亚群，并让真实使用者参与解释错误。评价报告还要说明未测试什么，防止通过一个醒目总分把证据边界抹掉。

\minitopic{升级路径比笼统警告更能保护使用者}“AI可能犯错，请核查”把全部负担交给用户，却不告诉他在何处、以何种来源核查。有效流程为不同风险设明确动作：低风险内容可自行复核，高风险冲突升级给领域专家，涉及权利或安全的决定自动暂停，并为紧急情况提供人工渠道。系统需把相关来源、冲突点和不确定类型一并传递，而不是只交出一个置信标签。升级是否成功还应以处理时间、撤回率和重复故障衡量。

\minitopic{供应商说明不能替代部署者验证}模型卡、数据说明和安全评估能提供重要元证据，却通常针对一般版本与测试环境。机构改变提示、接入本地数据库、设定自动执行权限或服务新群体后，风险结构已经变化。采购合同应要求版本通知、事件披露与审计支持，部署机构仍须验证自己的整合与工作流。责任不能因技术来自外部而外包，也不能因商业机密而使受影响者完全无法挑战结果。

\minitopic{停止使用也需要预先设计}组织容易建立上线审批，却没有谁能因错误积累而暂停系统。停用条件可包括关键亚群性能跌破底线、来源链失效、数据分布显著漂移、重大事故或监管要求变化；同时要准备人工替代流程，避免因业务依赖而让停用成为不可能。暂时停用不是承认技术永远失败，而是把修正能力置于速度之前。若系统只能不断更新却不能安全退出，它的制度可靠性本身就不足。

\minitopic{受影响者的申诉是一种认识输入}个体通常最早发现姓名错配、语境误读、罕见情况和现实后果，但机构若只把申诉看成客户服务，就会失去模型更新资料。申诉程序应允许查看决定依据、提交反证、获得人工复核并追踪修正；汇总数据则用于识别系统模式。个案陈述不是自动真实，也需核验，但其位置性经验能暴露测试集没有表示的故障。把申诉接入技术和政策更新，才使受影响者成为纠错参与者而非被动对象。

\begin{argumentbox}
可纠错制度形成一个闭环：部署前用真实任务验证；使用中保留来源与版本；异常时有升级和暂停；受影响者可以申诉；事故进入跨案例分析；修复后重新测试并公开适用边界。若只完成“上线前测试”，系统尚未具备持续值得信任的条件。
\end{argumentbox}

\minitopic{阶段性结论}可纠错的人机认识制度不追求让机器永不出错，而是让错误更易被发现、限制和学习。它公开证据链，按任务校准依赖，给审核者真实权限，检查分布性后果，记录版本和事故，并允许受影响者申诉。这样的制度把认识论从“输出像不像答案”推进到“整个系统是否值得信任、怎样知道自己何时不值得信任”。
